\section{Some norm inequalities}
\label{secframe}

In preparation for the proofs of Theorems \ref{th2} and \ref{th3} we
collect in this section a few useful lemmas.

\begin{lemma}
\label{Bbounded}
For $\alpha<-d$, the matrix
$\{b_{mn}^\alpha\}$ represents a bounded operator $B_\alpha$ on
$L^2(\mathbb R^d)$ with respect to the basis $\{|n\rangle\}$.
\end{lemma}
\begin{proof}
For any $x,y\in L^2(\mathbb R^d)$, consider their coordinate sequences
$\{x_n\},\{y_n\}\in \ell^2(\mathbb N_0^d)$ w.r.t.\ the basis $\{|n\rangle\}$. We have
\begin{gather*}
\vert\langle x,B_\alpha y\rangle\vert=
\bigg|\sum_{m,n}\overline{x_m} b^\alpha_{mn} y_n\bigg| \leq
\sum_{m,n}(1+|m-n|)^\alpha |x_m| |y_n|\\
\hphantom{\vert\langle x,B_\alpha y\rangle\vert=
\bigg|\sum_{m,n}\overline{x_m} b^\alpha_{mn} y_n\bigg|}{}
\leq\sum_{m,k}(1+|k|)^\alpha |x_m| |y_{k+m}|
\leq\sum_k(1+|k|)^\alpha \Vert x\Vert_{L^2} \Vert y\Vert_{L^2} ,
\end{gather*}
which concludes the proof.
\end{proof}

\begin{lemma}\label{anorm}
For any $\alpha>d$, there exists a constant $C_\alpha>0$ such that the following holds:
\[
\Vert\phi\Vert_a\leq C_\alpha \Vert\phi\Vert_{\infty,\alpha} ,\qquad \phi\in{\cal L}_{\infty,\alpha} .
\]
\end{lemma}

\begin{proof}
With notation as in the previous proof we have, for $\phi\in{\cal L}_{\infty,\alpha}$,
\begin{gather*}
\bigg|\sum_{m,n}|\phi_{mn}| \overline{x_n}\,y_m\bigg| \leq
\sum_{m,n}\big|b^\alpha_{mn} \phi_{mn}\big|
\big|b^{-\alpha}_{mn} \overline{x_n} y_m\big|
\leq \sup_{m,n}\big|b^\alpha_{mn}\phi_{mn}\big|
 \Vert B_{-\alpha}\Vert_{\rm op} \Vert x\Vert_{L^2} \Vert y\Vert_{L^2} ,
\end{gather*}
which evidently implies the claim by Lemma~\ref{Bbounded}.
\end{proof}

\begin{lemma}\label{trace}
For any $\alpha<-d$ there exists a constant $C'_\alpha>0$ such that
\[
\Vert\phi\Vert_{1,\alpha}\leq C'_{\alpha} \Vert\phi\Vert_1 ,\qquad \phi\in{\cal L}_1 ,
\]
where ${\cal L}_1$ denotes the space of trace class operators and
$\Vert\cdot\Vert_1$ is the standard trace-norm.
\end{lemma}

\begin{proof}
Writing the trace class operator $\phi$ as $\frac
12(\phi+\phi^*)+\frac i2(i\phi^*-i\phi)$ we may assume that $\phi$ is
self-adjoint. Then, writing $\phi= \phi_+-\phi_-$ where $\phi_\pm$ are positive
operators each of which has trace norm at most that of $\phi$, we can
assume $\phi$ is positive. In this case the Cauchy--Schwarz inequality
gives
\[
|\phi_{kl}|  \leq \big(\phi_{kk} \phi_{ll}\big)^{1/2},
\]
and hence
\begin{gather*}
\sum_{k,l}b_{kl}^{\alpha} |\phi_{kl}|\leq
\sum_{k,l}b_{kl}^{\alpha} \phi_{kk}^{1/2} \phi_{ll}^{1/2}
\leq \Vert B_{\alpha}\Vert_{\rm op} \sum_k\phi_{kk}
=\Vert B_{\alpha}\Vert_{\rm op} \Vert\phi\Vert_1 ,
\end{gather*}
which proves the claim thanks to Lemma~\ref{Bbounded}.
\end{proof}

\begin{lemma}\label{product}
$a)$  For any $\alpha\geq 0$ there exists a constant $C_{1,\alpha}>0$
such that
\begin{gather}\label{prodest1}
\Vert\phi \psi\Vert_{2,\alpha}\leq C_{1,\alpha}(\Vert\phi\Vert_{2,\alpha} \Vert\psi\Vert_a
+\Vert\phi\Vert_a \Vert\psi\Vert_{2,\alpha}),\qquad \phi, \psi\in {\cal L}_{2,\alpha} .
\end{gather}

$b)$  For any $\alpha> d$ there exists a constant $C_{2,\alpha}>0$
such that
\begin{gather}\label{prodest2}
\Vert\phi \psi\Vert_{1,\alpha}\leq C_{2,\alpha}(\Vert\phi\Vert_{2,2\alpha} \Vert\psi\Vert_2
+\Vert\phi\Vert_2 \Vert\psi\Vert_{2,2\alpha}),\qquad \phi,\psi\in {\cal L}_{2, 2\alpha} .
\end{gather}
\end{lemma}

\begin{proof}
First, note that for $\alpha\geq 0$
\begin{gather}
\label{Bsum}
b_{mn}^\alpha\leq c(\alpha)\,(b_{mk}^\alpha+b_{kn}^\alpha) ,
\qquad  m,n,k\in\mathbb N_0^d,
\end{gather}
where the constant $c(\alpha)$ depends only on $\alpha$.
We then have
\begin{gather*}
\Vert\phi \psi\Vert_{2,\alpha}^2 =\sum_{m,n}b^{2\alpha}_{mn}\bigg|\sum_k\phi_{mk}
\psi_{kn}\bigg|^2
 \leq
C\sum_{m,n}\bigg(\sum_kb^{\alpha}_{mk} |\phi_{mk}| \, |\psi_{kn}|
+\sum_kb^{\alpha}_{kn} |\phi_{mk}|\,|\psi_{kn}|\bigg)^2\\
\hphantom{\Vert\phi \psi\Vert_{2,\alpha}^2}{}
\leq
2C\sum_{m,n}\bigg(\sum_kb^{\alpha}_{mk} |\phi_{mk}|\,|\psi_{kn}|\bigg)^2
+2C\sum_{m,n}\bigg(\sum_kb^{\alpha}_{kn} |\phi_{mk}|\,|\psi_{kn}|\bigg)^2\\
\hphantom{\Vert\phi \psi\Vert_{2,\alpha}^2}{}
=2C\,\big(\Vert(U_{\alpha}\widetilde\phi) \widetilde\psi\Vert_2^2+
\Vert\widetilde\phi (U_{\alpha}\widetilde\psi)\Vert_2^2\big)\\
\hphantom{\Vert\phi \psi\Vert_{2,\alpha}^2}{}
\leq 2C \big(\Vert U_{\alpha}\widetilde\phi\Vert_2^2
\Vert\widetilde\psi\Vert_{\rm op}^2+
\Vert U_\alpha\widetilde\psi\Vert_2^2
\Vert\widetilde\phi\Vert_{\rm op}^2\big)
 \leq 2C \big(\Vert\phi\Vert_{2,\alpha}
\Vert\psi\Vert_a+\Vert\psi\Vert_{2,\alpha} \Vert\phi\Vert_a\big)^2 ,
\end{gather*}
where  $C= c(\alpha)^2$ and $U_\alpha:\mathcal H_\alpha\to\mathcal H_0$ is
the unitary operator def\/ined by (\ref{Ualpha}). This establishes~(\ref{prodest1}).

Using again~\eqref{Bsum}, we have
\begin{gather}
\Vert\phi \psi\Vert_{1,\alpha} =\sum_{m,n}b^\alpha_{mn}\bigg|\sum_k\phi_{mk}
\psi_{kn}\bigg| \leq
c(2\alpha)\sum_{m,n,k}b^{-\alpha}_{mn}\big(|b^{2\alpha}_{mk} \phi_{mk}|\,
|\psi_{kn}|+|\phi_{mk}|\,|b^{2\alpha}_{kn} \psi_{kn}|\big)\nonumber\\
\hphantom{\Vert\phi \psi\Vert_{1,\alpha}}{}
=c(2\alpha)\big(\Vert(U_{2\alpha}\widetilde\phi) \widetilde\psi\Vert_{1,-\alpha}+
\Vert\widetilde\phi (U_{2\alpha}\widetilde\psi)\Vert_{1,-\alpha}\big) .\label{34}
\end{gather}
 Lemma~\ref{trace} and a H\"older inequality now yield, for $\alpha>d$,
\[
\Vert(U_{2\alpha}\widetilde\phi) \widetilde\psi\Vert_{1,-\alpha}
\leq C'_{-\alpha}\Vert(U_{2\alpha}\widetilde\phi) \widetilde\psi\Vert_1
\leq C'_{-\alpha}\Vert U_{2\alpha}\widetilde\phi\Vert_2\Vert\psi\Vert_2
= C'_{-\alpha}\Vert\phi\Vert_{2,2\alpha}\Vert\psi\Vert_2 .
\]
Combining this with \eqref{34} gives \eqref{prodest2}.
\end{proof}

\section{Decay of free solutions}
\label{secth2}

The goal in this section is to prove Theorem \ref{th2}. The main
step in the proof is to establish appropriate time-decay estimates for the matrix elements
$\left(e^{-it{\bf\Delta}}\right)_{nm,kl}$ of the free time-evolution
operator with respect to the orthonormal basis
$\{(|n\rangle\langle m|)\}$ for ${\cal H}$. This is our f\/irst objective.

Since the Weyl map $W$ is unitary up to a constant factor by \eqref{Wunitary}, the matrix elements of the heat operator $e^{-t{\bf\Delta}}$, $t>0$, can be
obtained in closed form from the well-known  expression for
the heat kernel in Euclidean space and the relation \eqref{tria},
making use of the fact that the functions $W^{-1}(|n\rangle\langle
m|)$ can be computed explicitly in terms of Laguerre polynomials. The
result is given in \cite{Raimar} and reads, in case $d=1$,
\begin{gather*}
\left(e^{-t{\bf\Delta}}\right)_{nm,kl}=
\delta_{m+k,n+l}\sum_{v=0}^{\min\{m,l\}} C_{nm,kl,v}
 \frac{t^{m+l-2v}}{(1+t)^{m+k+d}} ,
\end{gather*}
where
\begin{gather}
\label{matconst}
C_{nm,kl,v}:=
\sqrt{{n\choose m-v}{k\choose l-v}
{m\choose m-v}{l\choose l-v}} ,
\end{gather}
for arbitrary non-negative integers $n$, $m$, $k$, $l$. By convention the
binomial coef\/f\/icient $n\choose m$ vanishes unless $0\leq m\leq n$. Note that the presence of
the Kronecker delta factor is a consequence of rotational invariance
of the heat kernel in two-dimensional space, which in the operator
formulation entails that $ e^{-t{\bf\Delta}}$ commutes with $e^{-isa^*a}$ for
$s\in \mathbb R$, where we have dropped the subscript on the
annihilation and creation operators when $d=1$.

Since ${\bf\Delta}$ is a positive,  self-adjoint operator we obtain
$\left(e^{-it{\bf\Delta}}\right)_{nm,kl}$ for $t\in\mathbb R$ by analytic
continuation in $t$, that is
\begin{gather}
\label{mateluni}
\left(e^{-it{\bf\Delta}}\right)_{nm,kl}=
\delta_{m+k,n+l}\sum_{v=0}^{\min\{m,l\}} C_{nm,kl,v}
 \frac{(it)^{m+l-2v}}{(1+it)^{m+k+d}} .
\end{gather}

We next observe that these  matrix
elements are expressible in terms of Jacobi polynomials.

\begin{lemma}
\label{jacobi}
For $d=1$ and $l\leq m,k\leq n$, we have
\begin{gather}\label{jacrep}
\left(e^{-it{\bf\Delta}}\right)_{nm,kl}=
\delta_{m+k,n+l}
\sqrt{\frac{n! l!}{m! k!}}
\frac{(it)^{m+l}}{(1+it)^{m+k+1}}\big(1+t^{-2}\big)^l
  P_{l}^{n-m,m-l}\left(\frac{t^2-1}{t^2+1}\right) ,
\end{gather}
where $P_l^{\alpha,\beta}(X)$, $l\in\mathbb N_0$, $\alpha,\beta\geq 0$,
$X\in[-1,1]$, are the classical Jacobi polynomials with standard
normalization~{\rm \cite{Szego}}.
\end{lemma}

\begin{proof}
For  $m+k=n+l$ and $n\geq m$ we have
\begin{gather*}
C_{nm,kl,v}
=\sqrt{\frac{n! l!}{m! k!}}{m\choose v}{l+n-m\choose l-v},
\end{gather*}
and consequently, for $l\leq m$,
\begin{gather}
\label{matelements}
\left(e^{-it{\bf\Delta}}\right)_{nm,kl}=
\delta_{m+k,n+l}\sqrt{\frac{n! l!}{m! k!}}
\frac{(it)^{m+l}}{(1+it)^{m+k+1}}
\sum_{v=0}^{l} {m\choose v}{l+n-m\choose l-v}
 \big(-t^{-2}\big)^{v} .
\end{gather}
Recall now \cite{Szego} that the classical Jacobi polynomials
$P_l^{\alpha,\beta}(X)$, $l\in\mathbb N_0,$ are orthogonal
w.r.t.\ the weight function
\begin{gather}\label{weight}
w(X)   =  (1-X)^\alpha(X+1)^\beta,\qquad X\in[-1,1],
\end{gather}
for f\/ixed values of $\alpha, \beta > -1$. For our purposes we may
restrict attention to integer values of $\alpha$ and  $\beta$. A  convenient
explicit form of $P_l^{\alpha,\beta}(X)$ with standard normalization is
\[
P_l^{\alpha,\beta}(X)=\sum_{j=0}^l{l+\alpha\choose l-j}{l+\beta\choose
  j}\left(\frac{X-1}{2}\right)^j\left(\frac{X+1}{2}\right)^{l-j} .
\]
With
\begin{gather}\label{X-t}
X=(t^2-1)/(t^2+1)\in [-1,1]
\end{gather}
this gives
\[
P_l^{\alpha,\beta}\left(\frac{t^2-1}{t^2+1}\right)
= \big(1+t^{-2}\big)^l\sum_{j=0}^l{l+\alpha\choose l-j}
{l+\beta\choose j}\big(-t^{-2}\big)^{j}
\]
and the result follows by comparison with~\eqref{matelements}.
\end{proof}

The representation~\eqref{jacrep} combined with rather recent
uniform estimates on the Jacobi polynomials may now be used to derive the
following bound on the matrix elements in question.

\begin{lemma}
\label{Uniest1}
For $d\geq 1$, there exists a constant $C_d$, independent of $n,m,k,l
\in\mathbb N_0^d$ and
$t\in\mathbb R$, such that
\begin{gather}\label{decay0}
\big|\left(e^{-it{\bf\Delta}}\right)_{nm,kl}\big| \leq  C_d |t|^{-\frac
d2} ,\qquad |t|\geq 1 .
\end{gather}
\end{lemma}

\begin{proof} From the def\/inition of ${\bf\Delta}$ it follows that the
  matrix elements in question factorize into those for the
  one-dimensional case, so it suf\/f\/ices to consider $d=1$.

 Since $e^{-it{\bf\Delta}}$ is unitary, it follows from
  \eqref{mateluni} that
 \[
\left(e^{-it{\bf\Delta}}\right)_{nm,kl} =  \left(e^{-it{\bf\Delta}}\right)_{kl,nm}.
\]
 Moreover, since
$W(\overline f)=W(f)^*$ and the heat kernel on $\mathbb R^2$ is
symmetric in its two arguments we likewise have
 \[
\left(e^{-it{\bf\Delta}}\right)_{nm,kl} =  \left(e^{-it{\bf\Delta}}\right)_{mn,lk} .
\]
These symmetries may also be checked directly from \eqref{mateluni} and \eqref{matconst}.

Taking into account that $m+k=n+l$ for non-vanishing matrix elements
we can therefore assume that $l\leq m,k\leq n$. Lemma~\ref{jacobi} then yields
\begin{gather}
\label{estimategene}
\left|\left(e^{-it{\bf\Delta}}\right)_{nm,kl}\right|
 = \delta_{m+k,n+l}
\sqrt{\frac{n! l!}{m! k!}}(1+t^2)^{-1/2}
\big(1+t^{-2}\big)^{(l-n)/2}|t|^{l-k}
\left|P_l^{\alpha,\beta}\left(\frac{t^2-1}{t^2+1}\right)\right|,
\end{gather}
where we have set
\begin{gather*}%\label{alphabeta}
\alpha:=n-m \qquad\mbox{and}\qquad \beta:=m-l .
\end{gather*}
Introducing the orthonormal Jacobi polynomials ${\bf
  P}_l^{\alpha,\beta}$ of degree $l\geq 0$ by normalizing w.r.t.\ the
$L^2$-norm def\/ined by the weight \eqref{weight}, one f\/inds (see e.g.~\cite[p.~67]{Szego})
\[
{\bf P}_l^{\alpha,\beta}(X)  =  \sqrt{\frac{(2l+\alpha +\beta +1)}{2^{\alpha
    +\beta +1}}\frac{(l+\alpha +\beta)! l!}{(l+\alpha)! (l+\beta)!}}
P_l^{\alpha,\beta}(X) ,
\]
and \eqref{estimategene} can be rewritten as
\begin{gather}
\left|\left(e^{-it{\bf\Delta}}\right)_{nm,kl}\right|
 = \delta_{m+k,n+l}\left(l+\frac{\alpha +\beta +1}2\right)^{-\frac 12}\nonumber\\
 \hphantom{\left|\left(e^{-it{\bf\Delta}}\right)_{nm,kl}\right|=}{}\times
\big(1+t^2\big)^{-1/2} (1-X)^{\alpha/2}(1+X)^{\beta/2}
\big|{\bf P}_l^{\alpha,\beta}(X)\big| ,\label{estimategene2}
\end{gather}
with $X$ given by \eqref{X-t}.

We now use the following uniform bound for the orthonormal Jacobi polynomials,
proven in~\cite{erdel}:
\begin{gather}\label{erdelyi}
(1-X)^{\alpha/2+1/4}
(1+X)^{\beta/2+1/4}\big|{\bf P}_l^{\alpha,\beta}(X)\big| \leq \sqrt{\frac{2 e}{\pi}}\sqrt{2+
\sqrt{\alpha^2+\beta^2}} ,
\end{gather}
valid for all $X\in]-1,1[$, $l\geq 0 $ and
$\alpha, \beta\geq-1/2$.

Applying this inequality in conjunction with \eqref{estimategene2} we
obtain, for $m,n,k,l\geq 0$ and   $|t|\geq 1$,
\begin{gather*}%\label{matest2}
\left|\left(e^{-it{\bf\Delta}}\right)_{mn,kl}\right|
 \leq C \delta_{m+k,n+l}
\left(\frac{2+
\sqrt{\alpha^2+\beta^2}}
{2l+\alpha +\beta +1}\right)^{\frac 12}
|t|^{-\frac 12}\leq C' \delta_{m+k,n+l} |t|^{-\frac 12} ,
\end{gather*}
which is the announced result.
\end{proof}
